\section{Multi-Scale Resolution Options}
\label{sec:options}

\subsection{Overview}

The multi-scale MCMC framework~\citep{autry2021multiscale} runs a ReCom chain
at a fine spatial level while using a coarser representation for
acceptance/rejection guidance.
This hierarchical structure reduces autocorrelation: coarse-level proposals
make larger jumps in plan space that would be rejected at full tract
resolution, but are tested at county resolution first and only expanded to
the fine level if they pass.

We define three multi-scale options corresponding to adjacent pairs in the
GEOID hierarchy.
Each option is specified by two CLI flags,
\texttt{--multiscale-fine} and \texttt{--multiscale-coarse},
combined with \texttt{--search multiscale}:

\begin{table}[h]
\centering
\caption{Multi-scale option summary}
\label{tab:options}
\begin{tabular}{clllll}
\toprule
Option & Fine & Coarse & Extra data & GEOID prefix & Output level & Status \\
\midrule
A & \bg{}   & \tract{}  & BG file    & $[:11]$      & BG           & Phase 2 \\
B & \tract{} & \county{} & None       & $[:5]$       & Tract        & Implemented \\
C & \bg{}   & \county{} & BG file    & $[:5]$       & BG           & Phase 2 \\
\bottomrule
\end{tabular}
\end{table}

\noindent
\textbf{Option B} is the uniquely data-free option: county adjacency is
derived entirely from the tract graph and GEOIDs, so no additional download
is required.
Options A and C require the block-group adjacency file
(\texttt{\{state\}\_bg\_adjacency\_\{year\}.pkl}),
which is fetched separately but from the same P.L.\ 94-171 source.

\subsection{Input Validation}

The CLI enforces that \texttt{--multiscale-fine} is strictly finer than
\texttt{--multiscale-coarse}.
The valid orderings are $\bg < \tract < \county$ (strictly finer
$=$ smaller average geographic area, larger GEOID prefix length).
Invalid orderings return an informative error:

\begin{verbatim}
  ERROR: --multiscale-fine county is not finer than
         --multiscale-coarse bg.
  Valid orderings (fine -> coarse): bg->tract,
         bg->county, tract->county.
\end{verbatim}

\noindent
The check is performed before any data loading, so no compute is wasted
on malformed invocations.

\subsection{Option B: Tract$\to$County (Implemented)}
\label{ssec:optionb}

Option B is invoked as:
\begin{verbatim}
  redist state --state TX --year 2020 --search multiscale \
    --multiscale-fine tract --multiscale-coarse county \
    --multiscale-steps 2000 --multiscale-alpha 0.3
\end{verbatim}

\paragraph{Data requirements.}
Only the standard tract adjacency file is required.
The county adjacency graph $\ec$ is built in-memory by \buildc{}
immediately after tract adjacency is loaded, in $O(|E_T|)$ time
(Section~\ref{sec:geoid}).
The county partition $\mathbf{f}$ is derived by \derive{}
with prefix length 5.

\paragraph{Multi-scale step sequence.}
At each chain step, the multi-scale algorithm proceeds as follows:

\begin{enumerate}
  \item Draw a ReCom proposal at the county level (using $\ec$ and
        county populations $\mathbf{p}_C$).
  \item If the county proposal is accepted (population-balanced at
        coarse tolerance \texttt{coarse\_tol}), expand it to the
        tract level: all tracts in the affected counties inherit
        the county assignment.
  \item Apply a fine-level ReCom perturbation within the expanded region
        to restore tract-level population balance.
\end{enumerate}

\paragraph{Coarse tolerance rationale.}
The coarse tolerance is set to $3\times$ the fine tolerance.
Formally, if the fine-level tolerance is $\varepsilon_f = 0.005$
(statutory 0.5\% one-person-one-vote bound), then the coarse tolerance
is $\varepsilon_c = 3\varepsilon_f = 0.015$.
The factor of 3 is justified as follows.
A county contains on average $n_T / n_C \approx 22$ tracts (for NC).
Population imbalance at the county level of $\varepsilon_c$ can be
corrected by the fine-level perturbation step without violating
$\varepsilon_f$, because the fine-level correction has 22 tracts of
granularity with which to rebalance.
A tighter coarse tolerance wastes accepted county proposals that could
be fine-corrected; a looser tolerance risks proposals the fine
perturbation cannot fix.
This is a practical heuristic, not a formal guarantee; the Phase 2
empirical study (Section~\ref{sec:empirical}) validates it on TX.

\paragraph{Output.}
The output plan is recorded at the fine (tract) level.
The manifest records \texttt{plan\_resolution = "tract"} and
\texttt{multiscale\_fine = "tract"}, \texttt{multiscale\_coarse = "county"}.

\subsection{Option A: BG$\to$Tract (Phase 2)}
\label{ssec:optiona}

Option A runs the fine chain at block-group resolution with tract-level
coarsening.
The BG adjacency file (\texttt{\{state\}\_bg\_adjacency\_\{year\}.pkl})
must be available.
The partition is derived by \derive{} with prefix length 11.

Option A is the highest-precision multi-scale mode: plans are output at
BG resolution ($\approx$5{,}000 units for NC), enabling finer district
boundaries than are possible at tract level.
This is particularly relevant for state legislative plans with many small
districts.

\textbf{Phase 2 status}: BG adjacency loading is implemented for the
single-resolution case (\texttt{--resolution block\_group});
the multi-scale wiring (passing BG adjacency as the fine graph to the
multi-scale chain) is planned for Phase 2.

\subsection{Option C: BG$\to$County (Phase 2)}
\label{ssec:optionc}

Option C combines the finest spatial precision (BG) with the coarsest
useful guidance level (county), skipping the intermediate tract level.
The BG-to-county partition is derived by \derive{} with prefix length 5
applied to BG GEOIDs (12 characters) --- the same prefix length as
tract$\to$county but applied to longer source strings.

Option C provides the most aggressive autocorrelation reduction (county
proposals make the largest jumps) while maintaining full BG-level output.
It requires both the BG adjacency file and BG population data.

\textbf{Phase 2 status}: Planned alongside Option A.

\subsection{Seeding and Reproducibility}

Multi-scale steps use the same \texttt{MSC\_STEP\_} seeding formula as
single-scale runs.
The \texttt{plan\_resolution} field is recorded in the manifest rather
than encoded in the seed: changing resolution with the same
\texttt{base\_seed} produces different plans (the adjacency graph differs),
which is correct behavior.
A user wishing to reproduce an Option~B run must supply the same
\texttt{base\_seed}, \texttt{--multiscale-steps}, and
\texttt{--multiscale-alpha} values alongside \texttt{--search multiscale
--multiscale-fine tract --multiscale-coarse county}.
All four parameters are recorded in the manifest.

\subsection{CLI Reference}

\begin{verbatim}
  # Tract level -- default, unchanged
  redist state --state NC --year 2020
  # Option B (no extra data)
  redist state --state TX --year 2020 --search multiscale \
    --multiscale-fine tract --multiscale-coarse county \
    --multiscale-steps 2000 --multiscale-alpha 0.3
  # Option A (BG file required)
  redist state --state NC --year 2020 --search multiscale \
    --multiscale-fine bg --multiscale-coarse tract
  # Option C (BG file required)
  redist state --state NC --year 2020 --search multiscale \
    --multiscale-fine bg --multiscale-coarse county
\end{verbatim}
